\begin{abstract}
Modern visual recognition requires both fine image detail and long-range
context, but retaining every patch through every transformer layer is costly
at high resolution. Token merging reduces this cost only when a learnable
merging policy also produces a physically shorter sequence. We introduce
\mn{}, a local-to-latent vision transformer that connects these two goals in
one forward path. A full-grid encoder first performs differentiable mass
routing over a learned similarity graph whose eligible edges are constrained
by the original two-dimensional patch geometry. A threshold-based soft
top-$k$ relaxation trains the routing and carrier scores, while an exact-budget
hard gather then materializes a shorter sequence before global reasoning.
Selected token mass is preserved as a logarithmic key bias in the latent
attention. The released operators also contain a rank-one query/key
augmentation that evaluates this bias with an unmodified FlashAttention
kernel; the reported ImageNet checkpoints use PyTorch fused attention for
this step, so we report the FlashAttention path separately.

On ImageNet-1K with DeiT-S/8, the trained \mn{} checkpoint reaches
81.368\% top-1 at $2.00\times$ patch compression; its matched-budget
re-evaluation is 81.360\%, compared with 82.246\% for dense DeiT-S/8 and
81.932\% for post-training ToMe. The exact matched DTEM score is
\pendingresult{DTEM top-1}. Radius-3 spatial routing improves the 150-epoch
ImageNet result by 1.670 percentage points over unconstrained routing and by
0.598 points over a flattened-window control. A separate three-seed
CIFAR-100 study reproduces the spatial advantage against a degree-matched
receiver permutation. Under a common trained-checkpoint
protocol, latency and peak-memory reductions remain
\pendingresult{latency} and \pendingresult{memory}. These experiments isolate
the benefit of a learned spatial routing prior while reserving deployment
claims until the missing checkpoint-side measurements are available.
\end{abstract}
